\documentclass[12pt]{article}
\usepackage[utf8]{inputenc}
\usepackage[T1]{fontenc}
\usepackage[margin=1in]{geometry}
\usepackage{setspace}
\usepackage{xcolor}
\usepackage[colorlinks=true, urlcolor=blue!60!black, linkcolor=blue!60!black]{hyperref}
\setlength{\parindent}{2em}
\setlength{\parskip}{1.25em}
\renewcommand{\baselinestretch}{1.0}
\begin{document}
\begin{center}
{\LARGE \textit{Moody v. NetChoice, LLC}, 603 U.S. 707 (2024)\par}
\vspace{0.8em}
{\large Case Brief\par}
\end{center}
\vspace{1em}

\section{Memory Jogger}
The Court unanimously vacated both lower-court judgments because neither court applied the demanding test for facial First Amendment challenges; a majority (Kagan, J.) added that curating third-party posts in a feed is protected editorial expression that a State may not override to rebalance the marketplace of ideas (\href{https://supreme.justia.com/cases/federal/us/603/22-277/}{Justia}).

\section{Detailed Case Facts}
In 2021, amid legislative concerns that large technology companies were suppressing conservative viewpoints, Florida and Texas enacted statutes restricting social-media content moderation and mandating individualized user explanations (\href{https://www.law.cornell.edu/supremecourt/text/22-277}{Cornell LII}):
\begin{itemize}
    \item \textbf{Florida S.B.~7072 (Fla.\ Stat.~\S~501.2041):} Covers ``social media platforms'' with over \$100 million in annual gross revenue or 100 million monthly active users. It bars platforms from ``censoring,'' deprioritizing, or labeling posts by or about political candidates or ``journalistic enterprises,'' requires ``consistent'' enforcement of content-moderation rules, limits rule changes to once every 30 days, mandates a chronological opt-out from sorting algorithms, and requires a written ``thorough rationale'' within seven days of any removal or alteration.
    \item \textbf{Texas H.B.~20 (Tex.\ Bus.\ \& Com.\ Code \S~120.001 \textit{et seq.}; Tex.\ Civ.\ Prac.\ \& Rem.\ Code \S~143A.001 \textit{et seq.}):} Covers platforms with over 50 million monthly U.S. users (excluding ISPs, email, and preselected news/entertainment sites). It forbids ``censoring''---blocking, banning, removing, deplatforming, demonetizing, or de-boosting---a user or user's expression based on viewpoint, and requires concurrent removal explanations plus a 14-day appeal process.
\end{itemize}

\section{Procedural History}
NetChoice, LLC and the Computer \& Communications Industry Association (collectively, NetChoice)---trade associations representing Facebook, YouTube, X, and others---brought pre-enforcement \emph{facial} First Amendment challenges in both States. Federal district courts in Florida (546 F. Supp.~3d 1082 (N.D. Fla.~2021)) and Texas (573 F. Supp.~3d 1092 (W.D. Tex.~2021)) preliminarily enjoined both laws as infringing platforms' protected editorial judgment. A circuit split ensued: the Eleventh Circuit largely affirmed Florida's injunction (34 F.4th 1196 (2022)), holding that content moderation is protected editorial discretion and that individualized explanations unduly burden speech under \textit{Zauderer v. Office of Disciplinary Counsel}, 471 U.S. 626 (1985), whereas the Fifth Circuit reversed Texas's injunction (49 F.4th 439 (2022)), declaring that platforms' content moderation is ``not speech'' at all and that Texas had a valid interest in ``protecting a diversity of ideas'' (\href{https://www.law.cornell.edu/supremecourt/text/22-277}{Cornell LII}).

\section{Judicial Votes}
The Supreme Court \textbf{vacated and remanded} the judgments of the Eleventh and Fifth Circuits by a \textbf{9--0} vote on the judgment. The case was argued on February 26, 2024, and decided on July 1, 2024 (\href{https://supreme.justia.com/cases/federal/us/603/22-277/}{Justia}; \href{https://www.law.cornell.edu/supremecourt/text/22-277}{Cornell LII}):

\begin{itemize}
    \item \textbf{Majority Opinion:} Justice Kagan delivered the opinion of the Court, joined in full by Chief Justice Roberts and Justices Sotomayor, Kavanaugh, and Barrett, and by Justice Jackson as to Parts I, II, and III--A.
    \item \textbf{Concurring Opinions:} Justice Barrett filed a concurring opinion. Justice Jackson filed an opinion concurring in part and concurring in the judgment. Justice Thomas filed an opinion concurring in the judgment. Justice Alito filed an opinion concurring in the judgment, joined by Justices Thomas and Gorsuch.
    \item \textbf{Dissenting Opinions:} None.
\end{itemize}

\section{Holding}
\textbf{Questions Presented:}

\begin{enumerate}
    \item \textbf{Facial Challenge Standard:} Did the Eleventh and Fifth Circuits properly evaluate NetChoice's \emph{facial} First Amendment challenges by focusing almost exclusively on Facebook's News Feed and YouTube's homepage without assessing the full range of platforms and functions covered by the statutes?
    \item \textbf{Protected Editorial Discretion:} Does a social-media platform's curation of third-party posts on its main feed constitute expressive activity protected by the First Amendment, and may a State restrict that curation to ``rebalance'' the marketplace of ideas?
\end{enumerate}

\textbf{Holding:}

\begin{enumerate}
    \item \textbf{No (Vacated and Remanded on Facial Standard):} Neither Circuit conducted the required facial First Amendment inquiry---determining each statute's full scope of covered platforms and functions and weighing whether a ``substantial number of [the law's] applications are unconstitutional, judged in relation to the statute's plainly legitimate sweep'' (\textit{Americans for Prosperity Foundation v. Bonta}, 594 U.S. 595, 615 (2021)) (\href{https://www.law.cornell.edu/supremecourt/text/22-277}{Cornell LII}).
    \item \textbf{Yes / No (Editorial Discretion Protected; ``Rebalancing'' Interest Invalid):} To guide remand and correct the Fifth Circuit's doctrinal errors, a majority held that curating third-party posts on curated feeds like Facebook's News Feed and YouTube's homepage is protected expressive activity under \textit{Tornillo}, \textit{PG\&E}, \textit{Turner}, and \textit{Hurley}, and that a State's asserted interest in correcting perceived ideological bias cannot survive even intermediate scrutiny (\href{https://www.law.cornell.edu/supremecourt/text/22-277}{Cornell LII}).
\end{enumerate}

\section{Analysis of Opinions}
\subsection{Majority Opinion (Justice Kagan)}

\subsection{The Demanding Framework for Facial First Amendment Challenges (Part II)}
Justice Kagan emphasized that facial challenges are disfavored because they rest on speculation and risk short-circuiting the democratic process (\textit{Washington State Grange v. Washington State Republican Party}, 552 U.S. 442, 450--51 (2008)). Although the First Amendment lowers the ordinary \textit{Salerno} ``no set of circumstances'' bar to provide breathing room for expression, a plaintiff must still show that the law ``prohibits a substantial amount of protected speech relative to its plainly legitimate sweep'' (\textit{United States v. Hansen}, 599 U.S. 762, 770 (2023)) (\href{https://www.law.cornell.edu/supremecourt/text/22-277}{Cornell LII}).

Both lower courts bypassed this two-step inquiry by treating the suits like as-applied challenges to Facebook's News Feed and YouTube's homepage. Step one requires defining the laws' full coverage---for instance, whether S.B.~7072 or H.B.~20 reaches direct messaging, Gmail, Etsy, Venmo, or Uber. Step two requires assessing First Amendment validity across each covered function (both for content moderation and for individualized explanations under \textit{Zauderer}) and comparing unconstitutional applications against constitutional ones. Because the Supreme Court is ``a court of review, not of first view'' and the record was underdeveloped, vacatur and remand were required.

\subsection{Editorial Discretion Over Curated Social-Media Feeds (Part III--A \& B)}
To prevent the Fifth Circuit from repeating its foundational First Amendment errors on remand, the Court synthesized its editorial-control precedents: \textit{Miami Herald Publishing Co. v. Tornillo}, 418 U.S. 241 (1974) (newspaper right-of-reply law struck down); \textit{Pacific Gas \& Electric Co. v. Public Utilities Comm'n of Cal.}, 475 U.S. 1 (1986) (utility billing-envelope insert mandate invalidated); \textit{Turner Broadcasting System, Inc. v. FCC}, 512 U.S. 622 (1994) (cable must-carry rules trigger First Amendment scrutiny); and \textit{Hurley v. Irish-American Gay, Lesbian and Bisexual Group of Boston, Inc.}, 515 U.S. 557 (1995) (parade organizers' exclusion of a contingent is protected expression). By contrast, \textit{PruneYard Shopping Center v. Robins}, 447 U.S. 74 (1980), and \textit{Rumsfeld v. FAIR}, 547 U.S. 47 (2006), involved hosts not engaged in their own expression (\href{https://www.law.cornell.edu/supremecourt/text/22-277}{Cornell LII}).

Three governing principles follow: (1) compiling, filtering, prioritizing, and labeling third-party speech into a curated product is itself protected expression; (2) a lenient host that admits most messages and excludes only a disfavored subset (as in \textit{Hurley}) does not forfeit First Amendment protection, nor does protection turn on a risk of misattribution; and (3) the government cannot override private editorial choices merely to ``better balance'' the marketplace of ideas. Applying these principles to Facebook's News Feed and YouTube's homepage, the Court explained that algorithms implementing Community Standards and Guidelines---removing, demoting, or labeling hate speech, self-harm, or election misinformation---produce a distinctive expressive offering that Texas's viewpoint-neutrality mandate directly alters. (The Court reserved feeds driven solely by user-engagement algorithms without independent content standards.)

\subsection{Invalidity of Texas's Asserted Interest (Part III--C)}
Without deciding whether strict or intermediate scrutiny governs, the Court held that Texas H.B.~20 fails even intermediate scrutiny under \textit{United States v. O'Brien}, 391 U.S. 367, 377 (1968), because Texas's stated objective---``to correct the mix of speech'' and stop platforms from ``silencing conservative viewpoints''---is directly related to the suppression of free expression. Distinguishing \textit{Turner} (which protected competition and local broadcast survival rather than altering speech mix), the Court reaffirmed that a State may not ``restrict the speech of some elements of our society in order to enhance the relative voice of others'' (\textit{Buckley v. Valeo}, 424 U.S. 1, 48--49 (1976)) (\href{https://www.law.cornell.edu/supremecourt/text/22-277}{Cornell LII}).

\subsection{Concurring Opinion (Justice Barrett)}

Justice Barrett joined the majority in full, warning that facial challenges force courts ``to bite off more than [they] can chew'' and urging as-applied suits. She flagged two critical open questions: (1) whether purely user-driven or autonomous AI/LLM moderation attenuates human editorial expression, and (2) whether foreign ownership and control abroad affect a platform's First Amendment rights under \textit{AID v. AOSI}, 591 U.S. 430 (2020) (\href{https://www.law.cornell.edu/supremecourt/text/22-277}{Cornell LII}).

\subsection{Concurring in Part and in the Judgment (Justice Jackson)}

Justice Jackson joined Parts I, II, and III--A, agreeing that the Eleventh Circuit stated First Amendment principles fairly while the Fifth Circuit did not, but declined to join Parts III--B, III--C, and IV under \textit{Ashwander} restraint rather than previewing merits rulings on an undeveloped record (\href{https://www.law.cornell.edu/supremecourt/text/22-277}{Cornell LII}).

\subsection{Concurring in the Judgment (Justice Thomas)}

Justice Thomas joined Justice Alito in full and wrote separately to urge reconsideration of \textit{Zauderer}, highlight the common-carrier doctrine for remand, and argue that \emph{all} facial challenges violate Article III's Case-or-Controversy limit by adjudicating nonparties' rights (\href{https://www.law.cornell.edu/supremecourt/text/22-277}{Cornell LII}).

\subsection{Concurring in the Judgment (Justice Alito)}

Justice Alito dismissed the majority's Part III merits analysis as ``nonbinding dicta'' and faulted NetChoice for failing to identify covered platforms, content types, or moderation algorithms. Arguing that \textit{PruneYard}, \textit{FAIR}, and common-carrier history deserve serious scrutiny, he questioned equating human newspaper editors with AI-driven platforms processing four petabytes of daily data (\href{https://www.law.cornell.edu/supremecourt/text/22-277}{Cornell LII}).

\subsection{Subsequent Developments}

\begin{itemize}

\item \textbf{Remand Litigation in Texas and Florida:} On remand, the Fifth Circuit sent Texas H.B.~20 back to the district court for ``thorough discovery'' across every statutory application (\textit{NetChoice, LLC v. Paxton}, 121 F.4th 494 (5th Cir.~2024)) (\href{https://www.supremecourt.gov/DocketPDF/25/25A97/368437/20250730143535035_No.%2025A97%20Response%20to%20Application.pdf}{SCOTUS Filing}). In Florida (\textit{NetChoice v. Uthmeier}), Judge Robert Hinkle denied Florida's motion to dismiss in May 2025 and denied cross-motions for summary judgment on June 19, 2026, advancing the case toward trial (\href{https://objectivelaw.org/issue-opportunities/moody-netchoice/}{Objective Law Initiative}).
    \item \textbf{Weaponization of \textit{Moody}'s Facial-Challenge Bar:} Lower courts have rigorously applied \textit{Moody} to vacate preliminary injunctions in trade-association facial challenges to state youth social-media and age-verification laws---including the Fifth Circuit in \textit{NetChoice, LLC v. Fitch}, 134 F.4th 799 (5th Cir.~2025) (Mississippi) (\href{https://www.supremecourt.gov/DocketPDF/25/25A97/368437/20250730143535035_No.%2025A97%20Response%20to%20Application.pdf}{SCOTUS Filing}) and the Eleventh Circuit in October 2026 vacating an injunction against Georgia's S.B.~351 for lacking \textit{Moody}'s full-scope comparison (\href{https://www.casemine.com/commentary/us/facial-challenges-to-social-media-age-restrictions-require-full-scope-analysis-under-moody/view}{CaseMine}).

\end{itemize}

\section{Examples: Future Applications}
\begin{enumerate}

\item \textbf{Same-Side (Facial Challenge Fails; As-Applied Claim Succeeds): A Viewpoint-Neutrality Mandate for ``All Large Online Platforms.''} A State enacts a statute barring ``all large online platforms'' from demoting or labeling user posts about climate policy based on viewpoint. A trade association brings a facial suit focusing solely on one member's curated newsfeed without addressing the statute's application to email providers, payment apps, or unmoderated message boards. Under \textit{Moody}, the court must reject facial invalidation until the challenger maps the statute's full scope and proves unconstitutional applications substantially outweigh constitutional ones; however, as applied to a curated newsfeed enforcing editorial guidelines, the statute violates the First Amendment because correcting ideological balance fails even intermediate scrutiny.

\item \textbf{Same-Side (Unconstitutional As Applied): A Must-Carry Rule for Candidates' Posts.} A State requires video-sharing platforms to host every post by a candidate for state office and forbids them to remove, demote, or label any such post. A video platform whose published guidelines bar election misinformation brings an \emph{as-applied} challenge after the State fines it for labeling a candidate's false claim about voting dates. Under \textit{Moody}, the platform's decisions about what to include, exclude, and label in its curated feed are protected editorial judgment, just as the newspaper's choices were in \textit{Tornillo}. Compelling it to carry and present candidates' posts unaltered changes its expressive product. The State's interest in giving candidates an equal voice is the kind of rebalancing interest \textit{Moody} held cannot justify overriding private editorial choices.

\item \textbf{Opposite-Side (Likely Valid): Political-Affiliation Protection for Payment Apps and Messaging Conduits.} Now suppose a State enacts a narrowly tailored statute applying \emph{solely} to private peer-to-peer payment apps and encrypted direct-messaging conduits (expressly excluding curated social-media feeds), prohibiting them from blocking financial transfers or private texts based on the sender's political affiliation, and requiring an automated one-line notice when a transaction is declined. An individual payment app brings an \emph{as-applied} challenge. Here, the outcome flips: (1) unlike Facebook's News Feed, a neutral payment conduit or direct-messaging pipe does not compile third-party speech into an expressive product (tracking \textit{PruneYard}, \textit{FAIR}, and \textit{Moody}'s own distinction between feeds and direct messaging); (2) regulating non-expressive conduit access invokes traditional common-carrier principles rather than \textit{Tornillo}/\textit{Hurley}; and (3) a brief, automated notice of a declined transfer is plausibly incidental to regulation of the transaction. If the notice separately compels commercial speech, the State must show that any applicable disclosure standard is satisfied, including \textit{Zauderer}'s limits on unjustified or unduly burdensome disclosures.

\item \textbf{Opposite-Side (Likely Valid): Protecting Ride-Share Drivers From Deactivation for Political Views.} A State law forbids ride-share companies to deactivate drivers because of political opinions they express off the app. A ride-share company sues, arguing under \textit{Moody} that the law interferes with its editorial judgment. The argument is weak. \textit{Moody}'s protection attaches to compiling and presenting third-party speech as a curated expressive product, such as a news feed. A ride-share service matches riders with drivers; deciding which drivers may use the app does not express any message, and the law does not require the company to display or endorse anyone's speech. As in \textit{PruneYard} and \textit{FAIR}, the regulated party is not engaged in its own expression, so the law is ordinary economic regulation outside \textit{Moody}'s editorial-discretion holding.

\item \textbf{Fence-Sitter (Open Question): A Feed Ranked Only by an Engagement Algorithm.} A State forbids platforms to rank posts on their main feeds based on the viewpoint expressed. A platform challenges the law as applied to a feed ranked entirely by an algorithm that maximizes predicted user engagement and applies no content standards or community guidelines of its own. For the platform: designing the algorithm is an editorial choice about how to present speech, and \textit{Moody} protects compiling and organizing third-party posts into a distinctive product. For the State: the Court expressly reserved feeds driven solely by user-engagement algorithms without independent content standards, and Justice Barrett questioned whether such automated sorting reflects any human expressive judgment. Because \textit{Moody} left the question open, the result is genuinely uncertain.

\end{enumerate}

\section{Critique}
\textbf{Holding vs.\ Dicta Tension:} Five Justices joined Part III's sweeping defense of algorithmic editorial discretion, yet three concurring Justices branded it ``nonbinding dicta'' and Justice Jackson abstained. Students must distinguish binding procedural vacatur from persuasive (or arguably necessary) remand guidance.

\textbf{The AI and ``Neutral Algorithm'' Gap:} Both the majority's footnote~6 and Justice Barrett's concurrence leave unresolved whether purely engagement-maximizing algorithms or autonomous generative-AI moderators qualify as human ``editorial judgment.''

\textbf{Procedural Asymmetry:} Requiring trade associations to map and litigate every marginal statutory application (from Gmail to Uber) before obtaining facial relief dramatically raises pre-enforcement litigation costs, encouraging targeted \emph{as-applied} suits by individual platforms instead.

\section{Key Quotations}
\begin{quote}
``In sum, there is much work to do below on both these cases, given the facial nature of NetChoice's challenges. But that work must be done consistent with the First Amendment, which does not go on leave when social media are involved.''
\end{quote}
\noindent --- Kagan, J. (\href{https://www.law.cornell.edu/supremecourt/text/22-277}{Cornell LII}).

\begin{quote}
``Deciding on the third-party speech that will be included in or excluded from a compilation---and then organizing and presenting the included items---is expressive activity of its own. And that activity results in a distinctive expressive product.''
\end{quote}
\noindent --- Kagan, J. (\href{https://www.law.cornell.edu/supremecourt/text/22-277}{Cornell LII}).

\begin{quote}
``On the spectrum of dangers to free expression, there are few greater than allowing the government to change the speech of private actors in order to achieve its own conception of speech nirvana.''
\end{quote}
\noindent --- Kagan, J. (\href{https://www.law.cornell.edu/supremecourt/text/22-277}{Cornell LII}).

\begin{quote}
``The holding in these cases is narrow: NetChoice failed to prove that the Florida and Texas laws they challenged are facially unconstitutional. Everything else in the opinion of the Court is nonbinding dicta.''
\end{quote}
\noindent --- Alito, J., concurring in the judgment (\href{https://www.law.cornell.edu/supremecourt/text/22-277}{Cornell LII}).


\end{document}
